% Generated draft: requires mathematical and source review.
\documentclass[11pt]{article}
\usepackage[margin=25mm]{geometry}
\usepackage{fontspec}
\setmainfont{Latin Modern Roman}
\usepackage{amsmath,amssymb,mathtools,unicode-math}
\setmathfont{Latin Modern Math}
\usepackage{enumitem,needspace,xurl,hyperref}
\setlength{\parindent}{0pt}
\setlength{\parskip}{5pt}
\setlength{\emergencystretch}{3em}
\newcommand{\N}{\mathbb{N}}
\newcommand{\Z}{\mathbb{Z}}
\newcommand{\Q}{\mathbb{Q}}
\newcommand{\R}{\mathbb{R}}
\newcommand{\C}{\mathbb{C}}
\newcommand{\F}{\mathbb{F}}
\newcommand{\A}{\mathbb{A}}
\newcommand{\PP}{\mathbb{P}}
\newcommand{\E}{\mathbb{E}}
\newcommand{\eps}{\varepsilon}
\newcommand{\dd}{\,\mathrm{d}}
\begin{document}
\section*{UnsolvedMath ID: 30002174}
\textbf{Source problem code:} catalogue entry
\subsection{Definitions and mathematical statement}
For a knot \(K\) and a color \(c\), write \(J_{K,c}(q)\) for the colored Jones polynomial and \(\mathcal H_{K,c}\) for the colored HOMFLY homology contemplated by the source. These symbols identify the intended invariants only: the source does not specify the class of knots, permitted colors, normalization, coefficient system, or grading conventions.

\textbf{Question as supplied.} Does the generalized or quantum volume conjecture for the colored Jones polynomial admit a natural categorification that yields a two-parameter super-\(A\)-polynomial governing the dependence of HOMFLY homology on color?

\textbf{Unresolved formulation.} The supplied record does not specify which generalized or quantum volume conjecture is intended, what constitutes a natural categorification, which two parameters are meant, or what mathematical relation expresses the proposed governing property. It also supplies no quantifier over knots or colors, no limiting regime, and no boundary or specialization conditions. Accordingly, this question records a proposed research direction; the evidence supplied does not determine a precise existence assertion, polynomial identity, recurrence, or asymptotic statement.
\subsection{Short English statement}
The question asks whether a version of the volume conjecture for colored Jones polynomials can be categorified so that a two-parameter super-\(A\)-polynomial controls how colored HOMFLY homology varies with color. The required constructions and governing relation remain unspecified, so the supplied record does not determine a precise mathematical proposition.
\subsection{Sources}
\begin{itemize}
\item Supplied local UnsolvedMath registry record, permanent ID \(30002174\), source problem code \texttt{catalogue entry}, title ``Categorification of the Quantum Volume Conjecture,'' in the Oberwolfach Reports Open Problems collection. The project identifies \texttt{data/unsolvedmath/problems.json} and the supplementary \texttt{catalog/problems.json} as local registries. No individual proposer is supplied; the proposed year is recorded as 2012.
\item Supplied report citation: \emph{Low-Dimensional Topology and Number Theory} (2013), Oberwolfach Reports, DOI \texttt{10.4171/owr/2012/42}, pp.\ 2541--2596. Supplied source URL and literature-source link labeled ``OWR report 42'': \url{https://doi.org/10.4171/owr/2012/42}. The background attributes the record to workshop open-problem discussions documented in Oberwolfach Reports; no specific talk author, passage, or problem page is identified.
\item Supplied literature assessment dated \texttt{2026-08-21}: the record is labeled open, but the assessment describes the extraction as a talk heading and broad categorification description rather than a precise proposition. It characterizes the generalized volume-categorification program as conjectural and the claim as uncheckable as stated. This is supplied dated curation, not an independently verified literature assessment; its time-sensitive status and cited claims require independent verification.
\end{itemize}
\end{document}
